\section{Implementation in \texttt{redist analyze}}
\label{sec:implementation}

\subsection{Command-Line Interface}

The three diagnostics described in this paper are implemented in the
\texttt{redist-analysis} crate of the \texttt{redist} Rust binary.
The command syntax is:

\begin{verbatim}
  redist analyze --types ensemble-diagnostics \
                 --state NC \
                 --year 2020 \
                 --chains 10 \
                 --steps 10000
\end{verbatim}

The command:
\begin{enumerate}
\item Loads the census tract adjacency graph for the specified state and year.
\item Generates $m = $ \texttt{--chains} random start plans via spanning-tree bisection.
\item Runs $n = $ \texttt{--steps} steps of \recom{} per chain.
\item Computes $\rhat$, $\ess$, and $\ham(1)$ at checkpoints.
\item Reports pass/fail against the thresholds from Section~\ref{sec:statutory}.
\end{enumerate}

\subsection{Output Format}

The output is a JSON file with the following structure:

\begin{verbatim}
{
  "state": "NC",
  "year": 2020,
  "k": 14,
  "chains": 10,
  "steps": 10000,
  "diagnostics": {
    "rhat_pp": 1.02,
    "rhat_seats": 1.03,
    "ess_pp": 769,
    "ess_seats": 1538,
    "hamming_lag1": 0.87
  },
  "thresholds": {
    "rhat_max": 1.1,
    "ess_min": 500,
    "hamming_max": 0.95
  },
  "verdict": "PASS"
}
\end{verbatim}

A \texttt{PASS} verdict indicates that the ensemble meets the evidentiary
standards of Section~\ref{sec:statutory} and is suitable for use in
redistricting litigation.

\subsection{Implementation Details}

\paragraph{ReCom proposal.}
The \recom{} step selects two adjacent districts uniformly at random,
merges their tract sets, generates a random spanning tree of the merged
subgraph, and bisects the spanning tree at its centroid edge subject to
population balance.
The implementation uses the \texttt{redist-core} spanning-tree sampler.

\paragraph{R-hat computation.}
Each chain's PP time series is stored in a ring buffer.
$\rhat$ is computed from the last half of each chain (disarding the
burn-in), following the recommendation of \citet{vehtari2021}.

\paragraph{ESS computation.}
$\ess$ uses the initial monotone sequence estimator of \citet{geyer1992}
to truncate the autocorrelation sum before it becomes noisy.

\paragraph{Hamming autocorrelation.}
The Hungarian algorithm for label canonicalization is $O(k^3)$ per plan
pair. For large $k$ we use the greedy approximation (sort districts by
centroid proximity), which is accurate to within 1\% for nearly-converged
chains.

\subsection{Relation to Ensemble Diagnostics Crate}

The diagnostics are exposed through the \texttt{ensemble\_diagnostics}
module in \texttt{redist-analysis}:

\begin{verbatim}
  redist_analysis::ensemble_diagnostics::{
      rhat,          // Gelman-Rubin R-hat
      ess,           // Effective sample size
      hamming_lag1,  // Lag-1 Hamming autocorrelation
      diagnose,      // Run all three and return verdict
  }
\end{verbatim}

The module has 21 unit tests covering boundary cases (all chains identical,
all chains diverged, single-chain degenerate input) and is part of the
standard \texttt{cargo test} suite.
